\section{Orthogonal components}
\label{sec:orthogonal}

In this section the prime is two.  We construct the coefficient required by
Lemma~\ref{lem:propagation} for orthogonal components.  The main choices are
the square classes of the norms in an orthogonal line decomposition.  They
control both the inner sign subgroup and the index lost on restriction.

\begin{proposition}\label{prop:orthogonal}
Let $G$ be a minimal-order counterexample to rational augmented reduced
homology nonvanishing at the prime two.  Thus $O_2(G)=1$ and
$\widetilde H_*(\mathcal A_2(G);\Q)=0$, whereas the asserted nonvanishing
holds for smaller groups with trivial $2$-core.  Then $G$ has no simple
component $L=\mathrm P\Omega(V)$, where $V$ is a nondegenerate symmetric
space over $\F_q$, the characteristic of $\F_q$ is at least five, and
its dimension $N$ is odd with $N\ge5$, even with $N\ge8$, or equal to six,
except possibly when $N=6$, $q=5$, and the determinant of the form is a
square.  In odd dimension $\mathrm P\Omega(V)=\Omega(V)$.
\end{proposition}

We use the determinant convention for discriminants: if $B$ is a Gram
matrix of the symmetric form, its discriminant is
$\det B\in\F_q^*/(\F_q^*)^2$.  We do not multiply this determinant by a
dimension-dependent sign.  In particular, for $N=2n$ the usual plus type
has determinant class $(-1)^n$, and minus type has the other class.  In
odd dimension we scale the form to make its determinant square.  This
does not change its isometry or similarity group, and it is possible
because scaling by $a$ multiplies the determinant by $a^N$.

Write $Q(v)=B(v,v)$ and let $\PO(V)$ denote the image of the isometry
group in $\PGL(V)$.  The spinor norm is normalized so that the reflection
in an anisotropic vector $v$ has spinor norm $Q(v)$ modulo squares.
The common kernel of determinant and spinor norm in the isometry group
is $\Omega(V)$, whose projective image is $L$; see
\citet[Lemma~11.49 and Theorems~11.50--11.51]{Taylor1992}.
These statements apply to the present isotropic spaces over odd
finite fields, all of dimension at least five.  We use the classical
automorphism theorem in its natural-module form: automorphisms of these
simple groups are induced by semilinear similarities, with the additional
diagram automorphisms in split dimension eight discussed below
\citep[Section~2 and statements~3.2--3.6]{Steinberg1960}; see also
\citet{GLS1998}. The groups in dimension six include the standard
isomorphisms
\[
 \mathrm P\Omega_6^+(q)\cong\PSL_4(q),\qquad
 \mathrm P\Omega_6^-(q)\cong\PSU_4(q).
\]
These identifications follow from the linear and unitary covers in
\citet[Corollaries~12.21 and~12.36]{Taylor1992}.

\subsection{Sign groups, parity, and their centralizers}

An orthogonal line frame is an unordered decomposition
$\mathcal D=\{\ell_1,\ldots,\ell_N\}$ into nondegenerate orthogonal
lines.  Color a line square or nonsquare according to the square class
of its nonzero norms.  Let $I$ be the set of nonsquare positions and put
$k=|I|$.  The group of independent signs on these lines, modulo their
common sign, is
\[
 S(\mathcal D)=\F_2^N/\langle j\rangle,
 \qquad j=(1,\ldots,1),
\]
embedded in $\PO(V)$.  We abbreviate its inner part to
$T(\mathcal D)=S(\mathcal D)\cap L$.  On $\F_2^N$ define
\[
 t(x)=\sum_{i=1}^N x_i,\qquad c(x)=\sum_{i\in I}x_i.
\]
The corresponding diagonal isometry has determinant $(-1)^{t(x)}$ and
spinor norm of square class $c(x)$.  Consequently the inverse image of
$T(\mathcal D)$ in $\F_2^N$ is
\begin{equation}\label{eq:orthogonal-inner-preimage}
 U=(\ker t\cap\ker c)+\langle j\rangle,
 \qquad
 U^\perp=\langle t,c\rangle\cap j^\perp.
\end{equation}
The addition of $\langle j\rangle$ accounts for the possibility of
replacing an isometry by its negative before testing membership in
$\Omega(V)$.

Here are the sign groups we shall use.  The homogeneous rows mean that
all lines have the same color; in the odd-dimensional row our determinant
normalization makes this the square color.  Mixed rows require both
colors to occur.
\begin{equation}\label{eq:orthogonal-rank-table}
\begin{array}{c|c|c|c}
\text{dimension and determinant}&\text{frame}&U^\perp&\rk T(\mathcal D)\\ \hline
N\text{ even, nonsquare}&k\text{ odd}&\langle t\rangle&N-2\\
N\text{ odd, square}&\text{homogeneous}&0&N-1\\
N\text{ even, square}&\text{homogeneous}&\langle t\rangle&N-2\\
N\text{ odd, square}&k\text{ even, mixed}&\langle c\rangle&N-2\\
N\text{ even, square}&k\text{ even, mixed}&\langle t,c\rangle&N-3
\end{array}
\end{equation}
Indeed $t(j)=N\bmod2$ and $c(j)=k\bmod2$, so the annihilators follow
directly from \eqref{eq:orthogonal-inner-preimage}; the ranks are
$\dim U-1$.  The discriminant of a frame is the product of its line
norm classes, so its being square is equivalent to $k$ being even.

We say that a projective sign subgroup distinguishes the positions if
no two coordinate characters agree on its inverse image in
$\F_2^N$.  Equality at positions $i\ne j$ would mean
$e_i^*+e_j^*\in U^\perp$.  Thus the first three rows of
\eqref{eq:orthogonal-rank-table} distinguish all positions for the
dimensions under consideration.  The fourth row does so when $k\ge4$.
The fifth does so when $k\ge4$ and $N-k\ge4$, because its three nonzero
annihilator vectors have supports $I$, $I^c$, and the entire set.

\begin{lemma}\label{lem:orthogonal-centralizer}
Let $T=T(\mathcal D)$ distinguish all $N$ positions and satisfy
$\rk T>N/2$.  Then every natural semilinear similarity centralizing $T$
projectively preserves each line of $\mathcal D$.  Its linear
centralizer is exactly $S(\mathcal D)$.  In dimension eight the same
conclusions hold for the full automorphism centralizer if some color
class has at least three members.  Every elementary abelian subgroup
of $\Aut(L)$ containing $T$ therefore lies in the natural semilinear
group, centralizes $S(\mathcal D)$, and has at most one direction beyond
its subgroup in $S(\mathcal D)$.
\end{lemma}

\begin{proof}
Let $\widetilde T$ be the inverse image of $T$ in the diagonal sign
group.  Its $N$ coordinate weights $\lambda_i$ are distinct.  Their
differences vanish on the common scalar sign and span $T^*$, so their
affine span has dimension $\rk T$.  A projective centralizer element
conjugates each sign matrix to itself times a scalar sign.  Field
automorphisms fix $1$ and $-1$.  Hence its permutation of the weight set
is translation by a single character $\mu\in T^*$.

If $\mu\ne0$, the $N$ weights fall into $N/2$ pairs
$\{\lambda,\lambda+\mu\}$.  Choosing one representative of each pair
shows that their affine span has dimension at most
$(N/2-1)+1=N/2$, a contradiction.  Thus $\mu=0$ and every weight line
is preserved.  A linear similarity preserving those lines is diagonal,
say with entries $a_i$.  Its common multiplier satisfies
$a_i^2=\eta$ for every $i$, and therefore $a_i/a_1\in\{1,-1\}$.
Its projective image belongs to $S(\mathcal D)$.

In split dimension eight, choose three lines in one color class.
The pair flips on the first two and on the last two belong to $T$,
since each has determinant one and square spinor norm.
Lemma~\ref{lem:orthogonal-eight-naturality} below shows that an
automorphism centralizing both is natural; it also gives naturality
of all automorphisms in minus dimension eight. In the other dimensions,
the natural-module automorphism theorem recalled above already gives
this conclusion. Thus the weight argument applies to every
automorphism centralizing $T$.

Finally, a semilinear map preserving every line centralizes all its
projective signs.  The field automorphism group is cyclic, so the image
of an elementary abelian $2$-group in it has rank at most one.  Its
linear kernel is contained in the sign group by the preceding
calculation.
\end{proof}

All applications below meet this lemma's rank hypotheses.  For the
first three rows of \eqref{eq:orthogonal-rank-table}, the minimum rank
is $N-2>N/2$ because $N\ge5$.  For the last row we use $N\ge8$, when
$N-3>N/2$.  The required color class in dimension eight always exists.

\subsection{Naturality in dimension eight}

The weight argument applies once an automorphism is known to act on
the natural space. In split dimension eight, two inner sign changes
supply that conclusion; in minus dimension eight all automorphisms
have natural realizations.

\begin{lemma}\label{lem:orthogonal-eight-naturality}
Let $V$ be a nondegenerate symmetric space of dimension eight over
$\F_q$, with defining characteristic at least five, and put
$L=\mathrm P\Omega(V)$. If $V$ has minus type, every automorphism of
$L$ is induced by a natural semilinear similarity of $V$.

If $V$ has plus type, let $\ell_1,\ell_2,\ell_3$ be mutually orthogonal
nondegenerate lines of the same norm square class. Let $x,y$ be the
projective isometries changing the signs on $\ell_1\oplus\ell_2$ and
$\ell_2\oplus\ell_3$, respectively, and fixing their orthogonal
complements. Then $x,y\in L$, and every automorphism of $L$ centralizing
both is induced by a natural semilinear similarity of $V$.
\end{lemma}

\begin{proof}
Suppose first that $V$ has plus type. We keep the two central quotients
of the spin group distinct. In split type $D_4$ the center of
$\Spin_8^+(q)$ has order four.  The kernel of its vector representation
\[
 \Spin_8^+(q)\longrightarrow\Omega_8^+(q)
\]
is the distinguished subgroup $\{1,-1\}$; the map onto
$L=\mathrm P\Omega_8^+(q)$ kills the full center.  Choose nonzero
vectors $v_i\in\ell_i$ for $1\leq i\leq3$. The pair flips $x,y$ belong
to $L$: each has determinant one and square spinor norm. In the Clifford
algebra their spin lifts are scalar multiples of $v_1v_2$ and $v_2v_3$. The scalar
normalizations can be chosen over $\F_q$ because the products of the
two norms are squares.  Orthogonality makes the vectors anticommute,
and these two lifts have commutator $-1$.

We now justify the covering and lifting assertions used in this
commutator test. The spin cover is the finite map
$\widetilde L=\Spin_8^+(q)\longrightarrow L=\mathrm P\Omega_8^+(q)$.
Its image and kernel can be checked without asserting surjectivity onto
all rational points of the adjoint algebraic group.  Reflection generation,
the parity of a reflection product, and the spinor kernel theorem
\cite[Theorems 11.39, 11.44, 11.50--11.51]{Taylor1992} show that
$\widetilde L\to \SO_8^+(q)$ has image $\Omega_8^+(q)$: an even reflection
product with square product of reflecting norms can be multiplied in the
Clifford algebra by a scalar in $\F_q$ to obtain norm one.
The vector action of an even Clifford normalizer has determinant one,
since it fixes the volume element.  Conversely, comparison with an even
reflection product shows that the vector image of a norm-one spin element
has square spinor norm.
Choose a hyperbolic basis $e_1,\ldots,e_4,f_1,\ldots,f_4$
in which $Q(\sum_i x_i e_i+y_i f_i)=\sum_{i=1}^4x_i y_i$, and put
\[
 z=\prod_{i=1}^4(e_i+f_i)(e_i-f_i).
\]
Then $z^2=1$, its Clifford norm is one, and its vector action is $-I$.
The vector kernel is the scalar subgroup $\{1,-1\}$, so the projective
kernel is $Z=\{1,-1,z,-z\}$; this is the full center of $\widetilde L$.

The group $\widetilde L$ is perfect. For a root $\alpha$, write
$x_\alpha(t)$ and $h_\alpha(a)$ for the unipotent and diagonal
elements in its root $\SL_2$. Root-subgroup generation and the calculation
\[
 h_\alpha(a)x_\alpha(t)h_\alpha(a)^{-1}x_\alpha(-t)
   =x_\alpha((a^2-1)t)
\]
put every root subgroup in $[\widetilde L,\widetilde L]$ once $a\in\F_q^\times$
satisfies $a^2\ne1$; such an $a$ exists for $q\ge5$
\cite[\S2, p.~416]{CEKT2013}.
Consequently an automorphism of $\widetilde L$ inducing the identity on $\widetilde L/Z$
has the form $s\mapsto s\chi(s)$ with a homomorphism $\chi:\widetilde L\to Z$,
and is the identity.  Thus any automorphism of $L$ has at most one lift
to $\widetilde L$; once the generator lifts below have been constructed, their
composites give a well-defined action on $Z$.

The completeness statement needed here is Steinberg's inner, diagonal,
field and graph decomposition, whose graph quotient in split type $D_4$
is $S_3$ \cite[\S3, statements~3.2--3.6]{Steinberg1960}.
The natural factors have lifts to $\widetilde L$.  Inner factors lift by conjugation.
For a similarity $g$ with multiplier $\mu$, the even-Clifford map
\[
 C_0(g)(uv)=\mu^{-1}g(u)g(v)
\]
lifts its projective conjugation action and fixes the scalar $-1$
\cite[(4.1), p.~422]{CEKT2013}.
This includes all diagonal factors. Let $\varepsilon_i$ be the
coordinate characters of the diagonal torus in the chosen hyperbolic
basis. For simple roots
$\varepsilon_1-\varepsilon_2$, $\varepsilon_2-\varepsilon_3$,
$\varepsilon_3-\varepsilon_4$, $\varepsilon_3+\varepsilon_4$ and any
prescribed factors $h_1,h_2,h_3,h_4\in\F_q^\times$, set
\[
 (a_1,a_2,a_3,a_4)=(h_1h_2h_3,h_2h_3,h_3,1),\qquad
 \mu=h_3/h_4.
\]
The proper similarity
$g=\operatorname{diag}(a_1,\ldots,a_4,\mu/a_1,\ldots,\mu/a_4)$
has exactly those root factors, since it scales the root parameters
for $\varepsilon_i-\varepsilon_j$ and $\varepsilon_i+\varepsilon_j$
by $a_i/a_j$ and $a_i a_j/\mu$, respectively.
Proper similarities fix $z$, as $\det(g)/\mu^4=1$; coefficient field
operations fix both $-1$ and $z$.  The improper isometry interchanging
$e_4$ and $f_4$ fixes $-1$, interchanges $z$ and $-z$, and interchanges
the two spin nodes of the diagram while fixing its vector node.

The remaining graph operation also lifts over $\F_q$.
The eight-dimensional para-Zorn algebra, with coordinates
$\alpha,\beta\in\F_q$ and $a,b\in\F_q^3$, has hyperbolic norm
$n(\alpha,a;b,\beta)=\alpha\beta-a\cdot b$ over the ground field
\cite[\S5, pp.~425--426]{CEKT2013}.  In the rational related-triple
description of $\widetilde L$, cyclic permutation of the three entries gives
triality and cycles the three nonidentity central sign triples
\cite[Proposition 4.6 and the following center description, p.~424]{CEKT2013}.
It preserves $Z$ and therefore descends to the finite quotient $\widetilde L/Z=L$.
Let $\mathcal N$ be the inner--diagonal--field subgroup in Steinberg's normal
sequence.  Its lifts fix $Z$ pointwise, whereas the descended triality
has order three and lies outside $\mathcal N$: otherwise uniqueness of lifts
would contradict its nontrivial action on $Z$.
The natural improper operation has a transposition action on
$Z\setminus\{1\}$ and also lies outside $\mathcal N$.
Their images generate $\operatorname{Aut}(L)/\mathcal N\cong S_3$, so they and
$\mathcal N$ generate $\operatorname{Aut}(L)$ and supply every automorphism with
a lift.  By uniqueness, the center action is a homomorphism that
identifies this quotient with $\operatorname{Sym}(Z\setminus\{1\})$.
The stabilizer of $-1$ is therefore precisely the natural subgroup:
it is the inverse image of the order-two subgroup represented by the
improper operation.  If an automorphism centralizes the two projective
pair flips, its lift sends each of their spin lifts to itself times an
element of $Z$.  It fixes their commutator $-1$, and hence is natural.

Now suppose that $V$ has minus type, and write $q=\ell^f$, where the
defining characteristic satisfies $\ell\ge5$.
In type ${}^2D_4$, the full-field factor in Steinberg's twisted construction
already includes the natural graph involution.  Explicitly, over
$E=\F_{q^2}$ use the split hyperbolic model and let $j$ interchange
$e_4,f_4$.  The fixed space
$W=\operatorname{Fix}(j\circ\operatorname{Frob}_q)$ has three hyperbolic
planes and the anisotropic plane with coordinates $(u,u^q)$ and norm
$u^{q+1}$, hence is an eight-dimensional minus space over $\F_q$.
The restriction $\tau=\operatorname{Frob}_{\ell}|_W$ is a natural semilinear
isometry of order $2f$, and $\tau^f=j|_W$ is the improper graph operation.
Therefore the full field group of $E$ in
\cite[\S2 and statement~3.5]{Steinberg1960} is realized naturally; the
absence of an additional graph factor for the twisted group omits no
natural graph automorphism.

The twisted diagonal factors are natural as well.  Their simple-root
factors have the form $h_1,h_2\in\F_q^\times$ and
$h_3=h$, $h_4=h^q$ with $h\in E^\times$.  Set
\[
 (a_1,a_2,a_3,a_4)=(h_1h_2,h_2,1,h^{-1}),\qquad
 \mu=(hh^q)^{-1}.
\]
The proper similarity
$g=\operatorname{diag}(a_1,\ldots,a_4,\mu/a_1,\ldots,\mu/a_4)$
has root factors $h_1,h_2,h,h^q$ by the same calculation as above.
Its first three hyperbolic pairs have coefficients in $\F_q$,
and its fourth pair is scaled by $(h^{-1},h^{-q})$.
Thus $g$ commutes with $j\circ\operatorname{Frob}_q$ and restricts to
a natural similarity of $W$, with multiplier $\mu\in\F_q^\times$.
Together with the inner and full-field factors, this realizes all
factors in Steinberg's twisted decomposition naturally.
\end{proof}

\subsection{Descent and the choice of a maximal configuration}

\begin{lemma}\label{lem:orthogonal-descent}
Suppose an involution $b$ of the natural projective semilinear group
has nontrivial field action and preserves every line of a nondegenerate
orthogonal frame of $V$.  Then $q=s^2$, and after scalar normalization
a representative of $b$ is an involutory semilinear map $B$ with fixed
space $V_0$ over $\F_s$.  A scalar multiple of the form descends to a
nondegenerate symmetric form on $V_0$.  The compatible frames are
exactly the scalar extensions of its nondegenerate orthogonal frames.
All their lines have one common norm square class over $\F_q$.
In particular, when $N$ is even the determinant over $\F_q$ is square.
\end{lemma}

\begin{proof}
The field action is the unique involution $\sigma:a\mapsto a^s$ of
$\F_q$.  A representative $B_1$ has $B_1^2=\lambda\id$; commuting
$B_1$ with its square gives $\lambda\in\F_s^*$.  The field norm
$\F_q^*\to\F_s^*$ is surjective, so multiplying $B_1$ by a scalar
makes its square one.  The fixed-vector descent statement in
Section~\ref{subsec:linear-compatible} then gives
$V=V_0\otimes_{\F_s}\F_q$ for $V_0=\{v:Bv=v\}$.

Write $Q(Bv)=\eta Q(v)^\sigma$.  The identity $B^2=1$ gives
$\eta\eta^\sigma=1$.  Hilbert's Theorem~90
\citep[Corollary~5.25]{MilneFields2022} supplies $a\in\F_q^*$
with $a^\sigma/a=\eta$.  Thus $Q_0=aQ$ satisfies
$Q_0(Bv)=Q_0(v)^\sigma$, and its restriction to $V_0$ is the descended
form.  It is nondegenerate because its scalar extension is
nondegenerate.  A $B$-stable line has a fixed nonzero vector by
one-dimensional descent.  This establishes the asserted correspondence
of frames, including nondegeneracy of their summands.

Every element of $\F_s^*$ is a square in $\F_q^*$, since $s+1$ is
even.  On a descended line the original form has a nonzero norm
$a^{-1}u$ with $u\in\F_s^*$; its ambient square class is consequently
the fixed class of $a^{-1}$.  Also
$\det Q=a^{-N}\det Q_0$.  For even $N$ both factors are squares over
$\F_q$, proving the last assertion.
\end{proof}

We now choose a family for Corollary~\ref{cor:maximal-family}.
Choose a nonempty set of allowed frames
whose inner signs meet Lemma~\ref{lem:orthogonal-centralizer}.  Consider
all elementary abelian $A\le N_G(L)$ such that $A\cap L$ contains
$T(\mathcal D)$ for some allowed frame and
\begin{equation}\label{eq:orthogonal-config}
 LA\text{ acts faithfully on }L,
 \qquad O_2(C_G(LA))=1.
\end{equation}
This family is nonempty: take $A=T(\mathcal D)\le L$, and use
Lemma~\ref{lem:component-core}. It is closed under enlargement to
faithful configurations, because an enlargement retains the contained
inner seed. Choose $A$ of maximum rank $r$, put
$H=LA$, and identify $H$ with its faithful image in $\Aut(L)$.

For its chosen frame, the centralizer lemma shows that $A$ preserves
every line and that its linear part is in $S=S(\mathcal D)$.  The group
$A$ contains $S\cap H$.  Indeed $S$ centralizes $A$, so adjoining any
element of $S\cap H$ preserves $H$, its centralizer, and the contained
inner seed; maximal rank forbids an enlargement.  Hence either
\begin{equation}\label{eq:orthogonal-A-no-field}
 A=S\cap H,
\end{equation}
or there is an involution $b\in A$ with nontrivial field action and
\begin{equation}\label{eq:orthogonal-A-field}
 E=S\times\langle b\rangle,\qquad
 A=E\cap H=(S\cap H)\times\langle b\rangle.
\end{equation}
In the latter case $H\cap\PO(V)$ is the entire linear part of $H$:
$H=L A$ and $A$ has linear kernel in $S\le\PO(V)$.

Our remaining task is to produce a cycle in $\mathcal A_2(H)$ with a
nonzero full-flag coefficient at some $A'$ of rank $r$ containing an
allowed seed and satisfying $LA'=H$.  Since $LA'=H$,
we have $O_2(C_G(LA'))=O_2(C_G(H))=1$, and the contained inner seed
gives $A'\cap L\ne1$. Thus $A'$ belongs to the same maximizing family.
Corollary~\ref{cor:maximal-family} then applies to this actual supported
subgroup.

For subsequent index calculations, let
\[
 P=\PO(V)/L,\qquad \pi:\PO(V)\longrightarrow P.
\]
Determinant and spinor norm identify this elementary abelian group with
\begin{equation}\label{eq:orthogonal-outer-signs}
 P\cong\F_2^2/\langle(N\bmod2,k\bmod2)\rangle.
\end{equation}
Both invariants are attained because in dimension at least five the
form represents both nonzero square classes.  A reflection in a square
line has image $(1,0)$, and one in a nonsquare line has image $(1,1)$.
Thus every mixed frame has $\pi(S)=P$.  For homogeneous frames of a
fixed color the image $\pi(S)$ is a fixed subgroup of $P$, independent
of the frame.  These observations will give constant indices except in
the small-field argument, where they give a maximum index.

\subsection{Nonsquare determinant and homogeneous frames}

First suppose $N$ is even and the determinant is nonsquare.  Every
orthogonal frame has an odd number of nonsquare lines and is mixed.
Use all these frames in the maximizing family.  Their inner signs have
rank $N-2$ and distinguish positions.  Lemma~\ref{lem:orthogonal-descent}
excludes a field axis, so \eqref{eq:orthogonal-A-no-field} holds and
$H\le\PO(V)$.  Put $B=H/L\le P$.  Every frame sign group surjects
onto $P$, whence
\[
 [S(\mathcal D):S(\mathcal D)\cap H]=[P:B],
 \qquad L(S(\mathcal D)\cap H)=H.
\]
The all-color case of Theorem~\ref{thm:frame-bound} supplies a nonzero
cycle of full flags on sign groups of rank $N-1$.  These are actual
coefficients: a full sign group recovers its frame as the common
eigenspaces of its diagonal lifts.  Lemma~\ref{lem:restriction}, at the
constant index $[P:B]$, gives the required cycle and coefficient at an
intersection $A'=S(\mathcal D)\cap H$.  The index exponent is at most
two and is smaller than $N-1$, as required.  All such intersections
have the original maximal rank and contain the allowed inner signs.
Corollary~\ref{cor:maximal-family} excludes this case. This argument
includes nonsquare determinant in dimension six for every $q$ under
consideration.

We may now assume that the determinant is square.  Suppose first that
$q\ge7$.  The classification of nondegenerate symmetric forms over a
finite field by dimension and determinant supplies homogeneous square
frames: the diagonal form with all entries one has the required
determinant.  Fix this color and maximize over configurations containing
its inner sign groups.  These have the second or third rank in
\eqref{eq:orthogonal-rank-table}.

If there is no field axis, use the one-color case of
Theorem~\ref{thm:frame-bound}.  To see the strict positivity at the
smallest parameter, an orthogonal plane generated by two lines of the
chosen color has
\[
 d_q=\frac{q-\chi_q(-1)}2\ge4
\]
 ordered frames of that color, where $\chi_q$ is the quadratic character
 of $\F_q^*$.  The one-color estimate has two roots
$\rho_+,\rho_-$ of $X^2-X+d_q^{-1}$, and its degree-$N$ fraction is
$\rho_+^N+\rho_-^N>0$.  This also holds when $d_q=4$, with the repeated
root $1/2$.  Thus a nonzero sign cycle exists.  If
$U=\pi(S(\mathcal D))$, then $U$ is constant over these frames and
$B=H/L\le U$ by the initial configuration.  Each intersection has
index $[U:B]$ and surjects onto $B$. The same coefficient detection and
constant-index restriction give a supported member of the original maximal rank, so
Corollary~\ref{cor:maximal-family} applies.

It remains in this case to treat a maximizing $A$ with field axis.
Fix its group $H$ and consider data $(\mathcal D,b)$, where
$\mathcal D$ is a homogeneous square frame and $b\in H$ is an
involution with the same nontrivial coset in
$H/(H\cap\PO(V))$, preserving its individual lines.  The initial
configuration supplies at least one such datum.  For a fixed $b$ which
occurs, Lemma~\ref{lem:orthogonal-descent} identifies all its compatible
frames with the orthogonal frames of a symmetric space over
$\F_s$, where $q=s^2$ and $s\ge5$.  Every such frame has the original
square color over $\F_q$; hence our restriction to that color discards
none of the descended frames.  We may apply the uniform all-color
fraction $(17/25)^N$ over $\F_s$, including both determinant types of
the descended space.

For a fixed frame, forgetting the axis has at least
\begin{equation}\label{eq:orthogonal-field-completions}
 |T(\mathcal D)|\ge 2^{N-2}
\end{equation}
completions.  Indeed all $bt$ with $t\in T(\mathcal D)$ are distinct
involutions in $H$ with the stipulated coset, centralize the signs, and
preserve the same frame.  Compatibility and its color are therefore
retained.  The needed strict comparison is
\begin{equation}\label{eq:orthogonal-field-inequality}
 \left(\frac{17}{25}\right)^N>2^{-(N-2)}\qquad(N\ge5).
\end{equation}
At $N=5$ the product of the left side with $2^{N-2}$ is greater than
one, and each increase in $N$ multiplies that product by $34/25>1$.

We verify detection before invoking the decoration count.  The group
$E=S(\mathcal D)\times\langle b\rangle$ recovers $S(\mathcal D)$ as
its kernel under the field map, and hence recovers $\mathcal D$.
In the split shuffle construction, flags beginning with
$\langle b\rangle$ and continuing with $\langle b\rangle U_i$, for a
sign flag $(U_i)$, recover the original sign coefficients.  A different
axis cannot contribute to such a flag with the same cone group: a
split rank-one subgroup outside the linear kernel is its own chosen
axis.  Thus all initial frame parameters are detected in actual chains.
Lemma~\ref{lem:decorations}, using
\eqref{eq:orthogonal-field-completions} and
\eqref{eq:orthogonal-field-inequality}, gives a nonzero homogeneous
cycle with full flags ending at groups $E$ of rank $N$.

Finally put $B=(H\cap\PO(V))/L$ and let $U$ be the fixed image of a
homogeneous square sign group in $P$.  The initial configuration shows
$B\le U$ and that its signs supply all of $B$.  For every new datum,
\[
 E\cap H=(S(\mathcal D)\cap H)\times\langle b\rangle,
 \quad [E:E\cap H]=[U:B],\quad L(E\cap H)=H.
\]
Restriction therefore has constant index, loses at most two ranks,
and returns an actual nonzero coefficient at a member of the
maximizing family of rank $r$. Corollary~\ref{cor:maximal-family} again
contradicts the choice of $G$.

\subsection{Forcing the required colors over the field of five elements}

We have reduced the proof of Proposition~\ref{prop:orthogonal} to
$q=5$ and square determinant.  There is no field direction.  For odd
$N\ge5$ call a frame admissible if its nonsquare count satisfies
\begin{equation}\label{eq:orthogonal-odd-colors}
 k\text{ even},\qquad 4\le k\le N-1.
\end{equation}
For even $N\ge8$ use
\begin{equation}\label{eq:orthogonal-even-colors}
 k\text{ even},\qquad 4\le k\le N-4.
\end{equation}
All these norm patterns occur by the classification by determinant.
Their inner signs distinguish positions and have rank $N-2$ in odd
dimension and $N-3$ in even dimension.  Consequently they meet
Lemma~\ref{lem:orthogonal-centralizer}.  Maximize configurations over
these admissible frames.  For the resulting $H$, its initial group is
$A=S\cap H$, and every admissible sign group surjects onto $P$.

The all-color frame theorem gives many cycles, but some of their
coefficients could be on inadmissible frames.  The following argument
shows that the produced coefficient space cannot be supported entirely
on those frames.  This is the remaining reason for the color ranges
\eqref{eq:orthogonal-odd-colors}--\eqref{eq:orthogonal-even-colors}.

\begin{lemma}\label{lem:orthogonal-five-colors}
Let $V$ be a square-determinant symmetric space over $\F_5$, of odd
dimension $N\ge5$ or even dimension $N\ge8$.  The space of actual sign
cycle coefficients produced by Theorem~\ref{thm:frame-bound} contains
an assignment with a nonzero full-flag coefficient at an admissible
frame as defined above.
\end{lemma}

\begin{proof}
Let $\mathcal F$ be the set of all orthogonal line frames and put
$M=|\mathcal F|$.  The produced coefficient space $\mathcal C$ has
dimension at least
\begin{equation}\label{eq:orthogonal-all-lower}
 \left(\frac{17}{25}\right)^N M(N-1)!.
\end{equation}
We use the description from Section~\ref{sec:frames}: local
coefficients are evaluations on binary merger brackets, with every
line letter odd.  A tree with a pair of leaf children satisfies the
plane relation obtained by replacing that pair by every orthogonal
frame of its span.  The brackets are precisely the evaluations detected
by the corresponding full partition flags.

Over $\F_5$ a square-determinant plane has one unordered frame of two
square lines and one of two nonsquare lines.  A nonsquare-determinant
plane has three unordered frames, each with one line of each color.
For completeness, use diagonal models $\langle1,1\rangle$ and
$\langle1,2\rangle$.  The projective lines of finite slope $t$ have
norms $1+t^2$ and $1+2t^2$, respectively, and the line of infinite slope
has norm $1$ or $2$.  The first plane has two isotropic lines, two
square lines, and two nonsquare lines; the second has three lines of
each nonzero norm type.  Orthogonal complementation pairs the lines
as asserted.  Since two odd letters satisfy $[x,y]=[y,x]$, the
length-two relation is the sum of one bracket for each unordered
replacement, with no cancellation between its two orderings.

Assume that all assignments in $\mathcal C$ vanish at every admissible
frame.  In odd dimension the remaining nonsquare counts are $0$ and
$2$; in even dimension they are $0,2,N-2,N$.  Call a remaining mixed
frame a two-minority frame: it has two lines of one color and $N-2$
of the other.  At such a frame, every coefficient of a bracket tree
whose bottom pair consists of two majority lines is zero.  Indeed the
other replacement in that same-color plane changes the majority pair
to the minority color and gives an admissible frame.  Its coefficient
is zero by assumption, and the two-term plane relation forces the
original coefficient to vanish.

We make the resulting dimension bound explicit.  Label the two
minority letters $a,b$ and the majority letters by a set $J$ of size
$N-2$.  The local functional factors through the multilinear part of
the free Lie superalgebra modulo
\[
 [u,v]=0\qquad(u,v\in J,\ u\ne v).
\]
Indeed the multilinear ideal generated by these relations is spanned
by bracket contexts containing a majority-majority bottom pair, whose
evaluations have just been shown to vanish.  For $u\in J$ set
$D_u=\operatorname{ad}(u)$.  The super Jacobi identity gives
\[
 D_uD_v+D_vD_u=\operatorname{ad}[u,v]=0.
\]
Thus the order of distinct majority adjoints on a branch matters only
by sign.  Fix a total order on $J$ and, for $I\subseteq J$, let $D_I$
denote their product in that order.  The local quotient is spanned by
\begin{equation}\label{eq:orthogonal-combs}
 [D_I(a),D_{J\setminus I}(b)]\qquad(I\subseteq J),
\end{equation}
so has dimension at most $2^{N-2}$.

To verify spanning, a subtree containing only majority letters and at
least two leaves is zero.  A subtree with exactly one minority letter
is therefore a comb rooted at that letter.  At the vertex where the
two minority branches first join, both branches are of this form.
Each majority letter above that vertex is distributed between the
branches by the super derivation identity.  Repeating and reordering
the adjoints gives \eqref{eq:orthogonal-combs}.  This also proves the
bound when additional local relations are present.

At a homogeneous frame choose any bottom pair in a merger tree.  Its
two-term plane relation expresses the coefficient in terms of the
coefficient at a two-minority frame.  Hence all assignments under our
vanishing assumption are determined by at most $2^{N-2}$ values per
two-minority frame.

We count those frames relative to $M$.  Every ordered pattern of norm
colors with an even number of nonsquare positions has the same number
of ordered frames: the isometry group acts transitively and the
stabilizer has order $2^N$, from independent signs on the lines.  There
are $2^{N-1}$ such patterns.  The fraction of unordered frames having
exactly two lines of a specified minority color is therefore
$\binom N2/2^{N-1}$.  There is one possible minority color in odd
dimension and two in even dimension.  Writing $e=1$ and $e=2$ in these
cases, respectively, we obtain
\begin{equation}\label{eq:orthogonal-excluded-upper}
 \frac{\dim\mathcal C}{M(N-1)!}
 \le
 e\frac{\binom N2}{2^{N-1}}
      \frac{2^{N-2}}{(N-1)!}
 =\frac{eN}{4(N-2)!}.
\end{equation}
For odd $N\ge7$ and even $N\ge8$ this contradicts
\eqref{eq:orthogonal-all-lower}.  At the two initial dimensions the
upper bounds are $7/480<(17/25)^7$ and $1/180<(17/25)^8$.
Increasing $N$ by two multiplies the upper bound by
$(N+2)/(N^2(N-1))$, which is smaller than $(17/25)^2$ throughout
these ranges.

For $N=5$ we use the mixed-pair plane relations as well.  At each
two-minority frame fix the numbering of its minority letters and the
ordering of its majority letters.  Introduce one formal coordinate
for each of the eight comb values in \eqref{eq:orthogonal-combs}.
Treating these coordinates as independent only enlarges the possible
coefficient space.  In every comb distribution at least one branch
contains a majority letter.  Reorder its adjoints so that this letter
is adjacent to the minority root.  Up to sign the coordinate is
therefore represented by a tree with a mixed bottom pair.

Its plane relation has three terms at three distinct two-minority
frames.  Each replacement again has a unique majority line and a
unique minority line in that plane, so the replaced tree is still
one of the two-comb distributions.  Reordering the majority adjoints
and, if necessary, interchanging the two minority branches converts
it to our chosen coordinate convention by a sign.  The relation thus
has three nonzero coefficients, each $1$ or $-1$, on three distinct
formal coordinates, including the selected coordinate.  In
particular, there is a relation with support at most three covering
each coordinate.

If the total number of formal coordinates is $D$, select a row basis
from these relations.  Its row supports still cover all $D$ columns:
a column zero on every basis row would be zero on every row.  Since
each basis row has support at most three, the rank is at least $D/3$.
The space of possible coordinate values has dimension at most $2D/3$.
Replacing \eqref{eq:orthogonal-excluded-upper} by this sharper bound
gives
\[
 \frac{\dim\mathcal C}{M\,4!}
 \le\frac23\frac5{24}=\frac5{36}
 <\left(\frac{17}{25}\right)^5,
\]
the final contradiction.  Thus some produced assignment has a
nonzero actual coefficient at an admissible frame.
\end{proof}

We finish the restriction and propagation in this last case.  The
lemma supplies a cycle on the full sign groups, of rank $m=N-1$, and
a nonzero coefficient at an admissible frame $\mathcal D$.  With
$B=H/L\le P$, every such frame has
\[
 [S(\mathcal D):S(\mathcal D)\cap H]=[P:B].
\]
An arbitrary frame $\mathcal D'$ occurring elsewhere in the same cycle
has $\pi(S(\mathcal D'))\le P$, and hence
\[
 [S(\mathcal D'):S(\mathcal D')\cap H]
 =[\pi(S(\mathcal D')):\pi(S(\mathcal D'))\cap B]
 \le[P:B].
\]
Our selected coefficient therefore has maximum intersection index
among all supported groups, even though the cycle also uses other
color counts.  Its index exponent is at most two, so it is strictly
less than $m$.  Lemma~\ref{lem:restriction} applies at this specific
coefficient and gives a cycle in $\mathcal A_2(H)$ with a nonzero full
flag ending at
\[
 A'=S(\mathcal D)\cap H.
\]
It contains the admissible inner seed, has the original maximal rank,
and satisfies $LA'=H$ because its signs supply all of $B$.  Restriction
has preserved an actual coefficient: in that lemma the fixed
complement makes intersection with $H$ injective on the residue's
flags. Corollary~\ref{cor:maximal-family} now applies to this very
supported subgroup $A'$ and gives the contradiction.

This proves Proposition~\ref{prop:orthogonal}.  The excluded case has
$N=6$, $q=5$, and square determinant, hence is plus type because
$-1$ is a square in $\F_5$.  Its simple group is $\PSL_4(5)$ and is
treated in Section~\ref{sec:exceptional}.
